\section{Discussion}

The account survives every test we can pose.
The site geometry, the wind climatology, the glide physics, the pitch stability, and
the impact biomechanics line up with what William reports, and they converge on the
same cause of failure that Eilmer himself named.
Some of these results hold regardless of our modelling choices and some depend on
them, and the two should not be read with equal confidence.

What holds firm does not depend on the fine detail of the apparatus.
The escarpment faces the prevailing wind and favourable launch conditions are common
(\cref{fig:site,fig:wind}); this is measured, not modelled.
The static instability of a tailless hand-and-foot wing is similarly hard to escape:
it follows from the pilot's CG lying aft of the wing's neutral point, which holds for
every posture and every plausible wing in the Monte Carlo
(\cref{fig:stability,fig:montecarlo}).
And the qualitative biomechanics are one-directional: a gentle landing cannot break
two legs, so the injury implies speed at impact (\cref{fig:biomech}).

What is model-dependent is the quantitative margin.
The ridge-lift magnitude rests on a first-order linear model over a smoothed
\SI{6}{\meter} DEM (\cref{eq:ridge}), so the \SI{1.5}{\meter\per\second} upwash and
the \SI{3.6}{\meter\per\second} wind threshold are order-of-magnitude statements that
a LIDAR-resolved large-eddy solve should sharpen.
The \SI{27}{\percent} reach fraction is a function of the priors we placed on wing
area, mass, and glide ratio, and should be read as the spread of a plausible space
rather than a calibrated probability.
The stability analysis is a linear short-period treatment of the pitch axis alone;
the lateral modes and the nonlinear post-departure motion need a full six-degree-of
freedom model before the \(75\times\) timescale separation can be called exact.

The Ibn Firnas parallel deserves a closer look in light of these results.
If two of the earliest credible glide attempts both ended in a tailless
loss-of-control arrival \autocite{anderson2024bridge}, the convergence is not
coincidence but physics: a tail is the one component an untutored builder is most
likely to omit and the one whose absence most reliably ends the flight.
The method here, terrain plus climate plus first-principles flight dynamics, can be
pointed at that account and others on the same footing.

The main threats to the conclusion are the unknown apparatus and the coarse terrain.
We mitigate the first with the Monte Carlo and the second by reporting slope and
lift as lower bounds, but neither is a substitute for the two obvious next steps: an
airborne-LIDAR terrain model with a microscale wind solve, and a structural analysis
of a willow-and-cloth wing under the computed flight loads to rule out a structural,
rather than control, failure.
We also do not reconstruct the actual flight.
We show that the reported flight is physically consistent and that its stated cause
is correct, which is a different and more defensible claim than asserting it happened
exactly as written.
